\renewcommand{\Anno}{2023}

\renewcommand{\NumeroProblema}{1}

\renewcommand{\Testo}{
Per ogni numero reale positivo $x$, indichiamo con $\operatorname{\mathsf{approx}}(x)$ la migliore approssimazione intera di $x$, definita in questo modo:
\begin{itemize}

\item  se nella rappresentazione decimale di $x$ la prima cifra dopo la virgola è 0, 1, 2, 3, 4, allora $\operatorname{\mathsf{approx}}(x)$ si ottiene da $x$ scartando la parte decimale,

\item  se nella rappresentazione decimale di $x$ la prima cifra dopo la virgola è 5, 6, 7, 8, 9, allora $\operatorname{\mathsf{approx}}(x)$ si ottiene da $x$ scartando la parte decimale e aumentando di uno la parte intera.

\end{itemize}

Così, per esempio, $\operatorname{\mathsf{approx}}(14)=\operatorname{\mathsf{approx}}(14{,}492)=14$ e $\operatorname{\mathsf{approx}}(33{,}5)=\operatorname{\mathsf{approx}}(33{,}502)=34$.


\begin{enumerate}
\renewcommand{\labelenumi}{(\alph{enumi})}

\item  Determinare tre numeri reali positivi $x_{1}$, $x_{2}$, $x_{3}$ tali che
$$x_{1}+x_{2}+x_{3}=1000
\qquad\quad\text{e}\quad\qquad
\operatorname{\mathsf{approx}}(x_{1})+\operatorname{\mathsf{approx}}(x_{2})+\operatorname{\mathsf{approx}}(x_{3})\neq 1000.$$

\item  Siano dati 500 numeri reali positivi $x_{1}$, $x_{2}$, \ldots, $x_{500}$ tali che
$$x_{1}+x_{2}+\ldots+x_{500}=1000.$$

Determinare il massimo ed il minimo valore possibile per
$$\operatorname{\mathsf{approx}}(x_{1})+\operatorname{\mathsf{approx}}(x_{2})+\ldots+\operatorname{\mathsf{approx}}(x_{500}).$$

\item  Determinare tutti i numeri reali positivi $x$ tali che
$$\frac{x}{2023}=\operatorname{\mathsf{approx}}\left(\sqrt{\operatorname{\mathsf{approx}}(x)}\right).$$

\end{enumerate}
}

\renewcommand{\Risposta}{
% SSSUP-BOZZA-AI: Risposta; autore=ChatGPT; data=2026-11-08; stato=DA-VERIFICARE
% SSSUP-PROVENIENZA: AI: elaborazione nuova dal testo
\begin{enumerate}
\item[(a)] Per esempio $333{,}4$, $333{,}4$, $333{,}2$: la somma degli arrotondamenti è $999$.
\item[(b)] Il minimo è $751$ e il massimo è $1250$.
\item[(c)] Le sole soluzioni sono $x=2023\cdot2022$ e $x=2023^2$.
\end{enumerate}
}

\renewcommand{\Aiutino}{
% SSSUP-BOZZA-AI: Aiutino; autore=ChatGPT; data=2026-11-08; stato=DA-VERIFICARE
% SSSUP-PROVENIENZA: AI: elaborazione nuova dal testo
La regola equivale a $\operatorname{\mathsf{approx}}(x)=\lfloor x+1/2\rfloor$. L'errore di arrotondamento appartiene a $(-1/2,1/2]$, con estremi diversi. Nel punto (c) il membro destro è intero, quindi $x$ deve essere un multiplo intero di $2023$.
}

\renewcommand{\SoluzioneUfficiale}{}

\renewcommand{\SoluzioniAlternative}{}

\renewcommand{\Osservazioni}{
% SSSUP-BOZZA-AI: Osservazioni; autore=ChatGPT; data=2026-11-13; stato=DA-VERIFICARE
% SSSUP-PROVENIENZA: rielaborazione redazionale AI da appunti di soluzione 2023, archivio privato Concorsi/SSSUP_2023/Soluzioni/SSSUP_2023_Sol.pdf, pagina 3
\TitoloOsservazione
Nel punto (b), gli estremi si possono realizzare anche con esempi particolarmente semplici, suggeriti negli appunti storici della prova.

Per il massimo $1250$, basta prendere
\[
x_1=\cdots=x_{499}=\frac12,\qquad x_{500}=750+\frac12.
\]
La somma dei $500$ numeri è $1000$, mentre gli arrotondamenti sono $499$ volte $1$ e una volta $751$, per un totale di $1250$.

Per il minimo $751$, fissiamo $0<\varepsilon<1/499$ e prendiamo
\[
x_1=\cdots=x_{499}=\frac12-\varepsilon,
\qquad x_{500}=750+\frac12+499\varepsilon.
\]
Anche in questo caso la somma è $1000$. I primi $499$ numeri si arrotondano a $0$, e l'ultimo a $751$. Questo realizza il minimo senza dover suddividere i numeri in più gruppi.
}

\renewcommand{\Commenti}{}

% Aggiunta redazionale: la soluzione preesistente non viene modificata.
\providecommand{\SoluzioneRedazionale}{}
\renewcommand{\SoluzioneRedazionale}{
% SSSUP-BOZZA-AI: SoluzioneRedazionale; autore=ChatGPT; data=2026-11-08; stato=DA-VERIFICARE
% SSSUP-PROVENIENZA: AI: elaborazione nuova dal testo
Scriviamo $R(x)=\lfloor x+1/2\rfloor$. Si ha sempre
\[
-\frac12<R(x)-x\leq\frac12.
\]

\paragraph{Punto (a).}
I tre numeri $333{,}4$, $333{,}4$, $333{,}2$ sommano $1000$, ma si arrotondano tutti a $333$.

\paragraph{Punto (b).}
La somma $M$ degli arrotondamenti è un intero e soddisfa $-250<M-1000\leq250$. Quindi $751\leq M\leq1250$. Il massimo si realizza prendendo $250$ numeri uguali a $1{,}5$ e $250$ uguali a $2{,}5$.

Per il minimo prendiamo $249$ numeri uguali a $1{,}498$ e $251$ uguali a $2{,}498$. La somma dei numeri è $249+502+500\cdot0{,}498=1000$, mentre la somma degli arrotondamenti è $249+502=751$.

\paragraph{Punto (c).}
Il membro destro dell'equazione è un intero non negativo; poiché $x>0$, possiamo scrivere $x=2023k$ con $k\geq1$ intero. Allora $R(x)=x$ e la condizione è
\[
k-\frac12\leq\sqrt{2023k}<k+\frac12.
\]
Elevando al quadrato si ottengono
\[
k(k-2024)+\frac14\leq0,\qquad
k(k-2022)+\frac14>0.
\]
Per interi positivi, la prima impone $k\leq2023$ e la seconda $k\geq2022$. Restano $k=2022,2023$, che soddisfano entrambe le condizioni.
}
